% ===========================================================================
% Paquets — jeu minimal couvrant les besoins réels du rapport.
% Toute addition doit être justifiée ; tout conflit se résout ici, jamais
% localement dans un chapitre.
% ===========================================================================

% ─── Langue, encodage, typographie ────────────────────────────────────────
% TODO(M-04) : confirmer la langue du rapport et les langues des résumés.
\usepackage[T1]{fontenc}
\usepackage[utf8]{inputenc}
\usepackage[french]{babel}
\usepackage{lmodern}
\usepackage{microtype}

% Latin Modern ne possède NI petites capitales grasses (`bx/sc`) NI petites
% capitales italiques (`m/scit`) — vérifié dans t1lmr.fd, qui ne déclare que
% `m/sc` et `m/scsl`. Sans les deux lignes ci-dessous, NFSS substitue quand
% même, mais prévient à chaque occurrence. On rend donc la substitution
% EXPLICITE, en conservant le rendu actuel : dans un en-tête de tableau le
% gras prime sur les petites capitales, et `scsl` est l'équivalent attendu de
% `scit`. Seul un changement de famille (newtxtext, …) donnerait de vraies
% petites capitales grasses ; ce n'est pas un choix à faire ici tant que le
% gabarit officiel n'impose pas sa police (TODO M-03).
% Déclaration différée : le fichier t1lmr.fd — donc la famille T1+lmr — n'est
% chargé qu'à la première sélection de police, qui a lieu après le préambule.
\AtBeginDocument{%
  \DeclareFontShape{T1}{lmr}{bx}{sc}{<->ssub*lmr/bx/n}{}%
  \DeclareFontShape{T1}{lmr}{m}{scit}{<->ssub*lmr/m/scsl}{}%
}

% ─── Géométrie et en-têtes ────────────────────────────────────────────────
% TODO(M-03) : marges imposées par le gabarit officiel.
\usepackage[a4paper,margin=2.5cm,headheight=15pt]{geometry}
\usepackage{fancyhdr}
\usepackage{setspace}

% ─── Mathématiques ────────────────────────────────────────────────────────
\usepackage{amsmath}
\usepackage{amssymb}

% ─── Graphiques, légendes, flottants ──────────────────────────────────────
\usepackage{graphicx}
\usepackage[font=small,labelfont=bf]{caption}
\usepackage{float}
\usepackage{rotating}          % figures en paysage (modèle de données)
\usepackage{adjustbox}         % filet de sécurité de largeur (cf. \pqfigure)

% ─── Diagrammes ───────────────────────────────────────────────────────────
% Les vues d'architecture sont décrites en TikZ plutôt qu'importées comme
% images : la source reste éditable et versionnée, et le rendu est vectoriel —
% double exigence de report/02_evidence_inventory.md.
% Bibliothèques : pointes de flèches, placement relatif, boîtes englobantes
% (frontières de confiance), calque d'arrière-plan, cylindres de stockage,
% entités du modèle de données, arithmétique de coordonnées.
\usepackage{tikz}
\usetikzlibrary{arrows.meta,positioning,fit,backgrounds,shapes.geometric,shapes.multipart,calc}

% ─── Tableaux ─────────────────────────────────────────────────────────────
\usepackage{array}
\usepackage{booktabs}
\usepackage{tabularx}
% `ragged2e` pour \RaggedRight. Différence décisive avec \raggedright : ce
% dernier met un ressort INFINI à droite, si bien que toute ligne « tient »
% et que TeX ne césure jamais — un mot plus long que la colonne déborde
% simplement en marge. \RaggedRight laisse un ressort fini, donc TeX coupe.
% C'est ce qui corrige les débordements de « Confidentialité »,
% « Implémentation », « Synchronisation »… dans les colonnes L étroites
% (cf. \newcolumntype{L} dans config/commands.tex).
\usepackage{ragged2e}
% NOTE — `longtable` est volontairement ABSENT.
% Sur l'installation de référence, longtable 4.27 (2026-01-28) est antérieur au
% noyau LaTeX 2026-06-01 et casse sur la terminaison de ligne : même
% `\begin{longtable}{llll} A & B \\ \hline` échoue. Le rapport n'emploie donc
% que `tabularx`, et les tables longues sont découpées par thème — ce qui se
% lit d'ailleurs mieux qu'une table de quatre-vingts lignes.
% Si la distribution est mise à jour (`miktex packages update`), longtable
% redevient utilisable ; rien dans le rapport n'en dépend.

% ─── Listes et extraits de code ───────────────────────────────────────────
\usepackage{enumitem}
\usepackage{listings}
\usepackage{multicol}          % listes de mots-clés sur deux colonnes

% ─── Couleur ──────────────────────────────────────────────────────────────
% Chargé explicitement : TikZ le tire déjà, mais la couverture et les
% marqueurs de champ à remplir s'en servent hors de tout environnement TikZ.
\usepackage{xcolor}

% ─── Acronymes ────────────────────────────────────────────────────────────
% NOTE — `glossaries` est volontairement ABSENT.
% Il était chargé sans jamais être employé : pas un \newacronym, pas un \gls
% dans tout le rapport. La liste d'acronymes de frontmatter/acronyms.tex est
% un `tabularx` écrit à la main, avec ses propres \section*.
% Trois raisons de le retirer plutôt que de le garder « au cas où » :
%   - il tirait datatool, translator et tracklang, d'où l'unique avertissement
%     résiduel de la compilation (« No `datatool' support for dialect
%     `french' ») ;
%   - il tirait aussi `longtable`, dont la version installée est cassée sur
%     cette distribution (voir la note de la section Tableaux ci-dessus) ;
%   - un paquet chargé sans usage est une dépendance qui ne se justifie pas.
% S'il faut un jour de VRAIS acronymes liés (renvois, première occurrence
% développée), le rétablir ici et convertir acronyms.tex en \newacronym.

% ─── Bibliographie ────────────────────────────────────────────────────────
% TODO(M-07) : style de citation imposé par l'établissement.
% `csquotes` AVANT biblatex, et non après : biblatex le détecte au chargement
% et l'utilise pour les guillemets des titres cités, qui suivent alors la
% convention française (« … ») fixée par babel.
\usepackage{csquotes}
\usepackage[backend=biber,style=numeric,sorting=none]{biblatex}

% ─── Hyperliens : à charger en dernier ────────────────────────────────────
\usepackage[hidelinks]{hyperref}
